\section{推断因果关系的局部准则}

IC* 算法以一个分布 $\hat{P}$ 为输入，输出一个部分有向图。其中有些连线被标记为单向（表示真正的因果），有些是不带标记的单向（表示可能的因果），有些是双向的（表示虚假关联），还有些是无向的（表示仍未能确定的关系）。产生这些标记的条件，可以被取作各类因果关系的定义。本节给出由 IC* 算法浮现出来的“可能因果”与“真正因果”的明确定义。注意，在所有这些定义中，判定两变量（$X$ 与 $Y$）之间因果关系的准则，都要求第三个变量 $Z$ 与 $X$、$Y$ 呈现某种特定的依赖模式。这并不奇怪，因为因果断言的实质就在于规定 $X$ 与 $Y$ 在第三个变量影响之下的行为，而该变量对应于对 $X$（或 $Y$）的外部控制——这正呼应了“没有操纵就没有因果”的范式（Holland 1986）。区别只在于：充当虚拟控制的变量 $Z$ 必须在数据本身之中被识别出来，仿佛自然已经做了那个实验。IC* 算法可以被看作提供了一种系统的方式，在稳定性假设下搜寻有资格充当虚拟控制的变量 $Z$。

\begin{definition}[可能原因]
若下列条件成立，则称变量 $X$ 对另一变量 $Y$ 有\emph{可能的因果影响}（可由 $\hat{P}$ 推断）：
\begin{enumerate}[label=\arabic*.,leftmargin=1.8em]
\item $X$ 与 $Y$ 在每个语境中都相依；
\item 存在变量 $Z$ 与语境 $S$，使得\\
(i) $X$ 与 $Z$ 在给定 $S$ 时独立（即 $(X\indep Z\mid S)$），且\\
(ii) $Z$ 与 $Y$ 在给定 $S$ 时相依（即 $(Z\nindep Y\mid S)$）。
\end{enumerate}
\end{definition}

所谓“语境”，我们指一组被赋予具体取值的变量。例如在图 2.3(a) 中，变量 $b$ 之所以有资格作为 $d$ 的可能原因，是因为变量 $Z=c$ 在语境 $S=a$ 中既与 $d$ 相依又与 $b$ 独立。同样，$c$ 也有资格作为 $d$ 的可能原因（取 $Z=b$、$S=a$）。$b$ 与 $c$ 都没有资格作为 $d$ 的\emph{真正}原因，因为这一依赖模式也与潜共同原因相容，即图 2.4(a)—(b) 中的双向弧。然而，定义 2.7.1 取消了 $d$ 作为 $b$（或 $c$）之原因的资格，这就导致把 $d$ 归类为 $e$ 的真正原因，如定义 2.7.2 所述。\footnote{定义 2.7.1 在 Pearl (1990) 中是作为事件（而非变量）之间的关系表述的，并附加了条件 $P(Y\mid X)>P(Y)$（本于 Reichenbach 1956、Good 1961 与 Suppes 1970 的精神）。这一精炼适用于本节的任何定义，但将不再显式表述。}注意，定义 2.7.1 排除了 $X$ 成为自身之可能原因、或成为在功能上决定 $X$ 的任何其他变量之可能原因。

\begin{definition}[真正原因]
若存在变量 $Z$ 使得下列之一成立，则称变量 $X$ 对另一变量 $Y$ 有\emph{真正的因果影响}：
\begin{enumerate}[label=\arabic*.,leftmargin=1.8em]
\item $X$ 与 $Y$ 在任何语境中都相依，且存在满足下列条件的语境 $S$：\\
(i) $Z$ 是 $X$ 的可能原因（依定义 2.7.1），\\
(ii) $Z$ 与 $Y$ 在给定 $S$ 时相依（即 $(Z\nindep Y\mid S)$），且\\
(iii) $Z$ 与 $Y$ 在给定 $S\cup X$ 时独立（即 $(Z\indep Y\mid S\cup X)$）；
\item $X$ 与 $Y$ 处于准则 1 所定义关系的传递闭包之中。
\end{enumerate}
\end{definition}

条件 (i)—(iii) 在图 2.3(a) 中以 $X=d$、$Y=e$、$Z=b$、$S=\varnothing$ 得到例示。通过以 $d$ 为条件而破坏 $b$ 与 $e$ 之间的依赖，这一事实无法归因于 $d$ 与 $e$ 之间的虚假关联；真正的因果影响是唯一的解释，如图 2.4 的各结构所示。

\begin{definition}[虚假关联]
若两变量 $X$ 与 $Y$ 在某个语境中相依，且存在另外两个变量（$Z_1$ 与 $Z_2$）与两个语境（$S_1$ 与 $S_2$）使得：
\begin{enumerate}[label=\arabic*.,leftmargin=1.8em]
\item $Z_1$ 与 $X$ 在给定 $S_1$ 时相依（即 $(Z_1\nindep X\mid S_1)$）；
\item $Z_1$ 与 $Y$ 在给定 $S_1$ 时独立（即 $(Z_1\indep Y\mid S_1)$）；
\item $Z_2$ 与 $Y$ 在给定 $S_2$ 时相依（即 $(Z_2\nindep Y\mid S_2)$）；且
\item $Z_2$ 与 $X$ 在给定 $S_2$ 时独立（即 $(Z_2\indep X\mid S_2)$）；
\end{enumerate}
则称 $X$ 与 $Y$ 之间是\emph{虚假关联}。
\end{definition}

条件 1 与 2 用 $Z_1$、$S_1$ 取消 $Y$ 作为 $X$ 之原因的资格，与定义 2.7.1 的条件 (i)—(ii) 平行；条件 3 与 4 用 $Z_2$、$S_2$ 取消 $X$ 作为 $Y$ 之原因的资格。于是，潜共同原因的存在便成为 $X$ 与 $Y$ 之间观测依赖的唯一解释，如结构 $Z_1\to X\leftarrow Y\leftarrow Z_2$ 所示。

当有时间信息可用时（大多数概率因果理论都作此假定——Suppes 1970; Spohn 1983; Granger 1988），定义 2.7.2 与 2.7.3 会大大简化，因为凡在 $X$ 之前且与 $X$ 邻接的变量，如今都有资格作为 $X$ 的“可能原因”。此外，只要语境 $S$ 被限定在 $X$ 之前，就不再要求邻接（即定义 2.7.1 的条件 1）。这些考虑导致区分真正原因与虚假原因的更简单条件，如下所示。

\begin{definition}[带时间信息的真正因果]
若存在第三个变量 $Z$ 与语境 $S$，二者都出现于 $X$ 之前，且满足：
\begin{enumerate}[label=\arabic*.,leftmargin=1.8em]
\item $(Z\nindep Y\mid S)$；
\item $(Z\indep Y\mid S\cup X)$；
\end{enumerate}
则称变量 $X$ 对 $Y$ 有\emph{因果影响}。
\end{definition}

定义 2.7.4 背后的直觉与定义 2.7.2 相同，只是现在用时间在先来确定 $Z$ 是 $X$ 的可能原因。这一点在图 2.5(a) 中得到例示：若以 $X$ 为条件能使 $Z$ 与 $Y$ 由相依转为独立（在语境 $S$ 中），那么 $Z$ 与 $Y$ 之间的依赖必定是以 $X$ 为中介的；既然 $Z$ 先于 $X$，这种中介便意味着 $X$ 对 $Y$ 有因果影响。

\begin{definition}[带时间信息的虚假关联]
若两变量 $X$ 与 $Y$ 在某个语境 $S$ 中相依，$X$ 先于 $Y$，且存在满足下列条件的变量 $Z$：
\begin{enumerate}[label=\arabic*.,leftmargin=1.8em]
\item $(Z\nindep Y\mid S)$；
\item $(Z\indep X\mid S)$；
\end{enumerate}
则称 $X$ 与 $Y$ 之间是\emph{虚假关联}。
\end{definition}

图 2.5(b) 例示了定义 2.7.5 背后的直觉。这里 $X$ 与 $Y$ 之间的依赖不能归因于二者之间的因果联系，因为这种联系将蕴含 $Z$ 与 $Y$ 之间的依赖，而条件 1 已将其排除。\footnote{回想因果依赖的传递性由稳定性所蕴含。虽然可以构造出 $Z$ 与 $Y$ 独立的因果链 $Z\to X\to Y$，但这种独立不会对链的一切参数化都持续成立。}

\begin{figure}[htbp]
\centering
\begin{tikzpicture}[causal,scale=0.95,every node/.style={transform shape}]
  % (a)
  \begin{scope}
    \node[circle,fill,inner sep=1.3pt] (Z) at (-1.3,1.15) {}; \node[font=\small,left=0.4mm] at (-1.3,1.15) {$Z$};
    \node[circle,fill,inner sep=1.3pt] (S) at (1.3,1.15) {};  \node[font=\small,right=0.4mm] at (1.3,1.15) {$S$};
    \node[circle,fill,inner sep=1.3pt] (X) at (0,0.15) {};    \node[font=\small,right=0.5mm] at (0,0.15) {$X$};
    \node[circle,fill,inner sep=1.3pt] (Y) at (0,-0.95) {};   \node[font=\small,below=0.3mm] at (0,-0.95) {$Y$};
    \draw[ed] (Z)--(S); \draw[ed] (Z)--(X); \draw[ed] (S)--(X);
    \draw[ed] (S)--(Y); \draw[ed] (X)--(Y);
    \node[font=\small] at (0,-1.7) {(a)};
  \end{scope}
  % (b)
  \begin{scope}[xshift=6.6cm]
    \node[circle,fill,inner sep=1.3pt] (S2) at (0,1.75) {};  \node[font=\small,above=0.3mm] at (0,1.75) {$S$};
    \node[circle,fill,inner sep=1.3pt] (Z2) at (-1.2,1.0) {}; \node[font=\small,left=0.4mm] at (-1.2,1.0) {$Z$};
    \node[circle,fill,inner sep=1.3pt] (W2) at (1.2,1.0) {};  \node[font=\small,right=0.4mm] at (1.2,1.0) {$W$};
    \node[circle,fill,inner sep=1.3pt] (X2) at (0,0.1) {};    \node[font=\small,right=0.5mm] at (0,0.1) {$X$};
    \node[circle,fill,inner sep=1.3pt] (Y2) at (0,-0.95) {};  \node[font=\small,below=0.3mm] at (0,-0.95) {$Y$};
    \draw[ed] (S2)--(Z2); \draw[ed] (S2)--(W2);
    \draw[ed] (Z2)--(X2); \draw[ed] (W2)--(X2);
    \draw[ed] (Y2)--(X2);
    \node[font=\small] at (0,-1.7) {(b)};
  \end{scope}
\end{tikzpicture}
\caption{图示时间信息如何使人们得以从 (a)、(b) 所示的条件独立中分别推断出（$X$ 与 $Y$ 之间的）真正因果与虚假关联。}
\label{fig:2.5}
\end{figure}

审视刚才给出的定义可见，一切因果关系都至少由三个变量推断而来。具体而言，使我们得以断定一个变量不是另一变量之因果后果的信息，以“非传递三元组”的形式出现——例如图 2.1(a) 中满足 $(a\indep b\mid\varnothing)$ 与 $(a\nindep c\mid\varnothing)$、$(b\nindep c\mid\varnothing)$ 的变量 $a,b,c$。论证如下。若我们找到条件（$S_{ab}$），使变量 $a$ 与 $b$ 各自与第三个变量 $c$ 相关、却彼此独立，那么第三个变量不可能是 $a$ 或 $b$ 的原因（回忆：在稳定分布中，共同原因的存在蕴含结果之间相依）；相反，$c$ 或者是它们的共同结果（$a\to c\leftarrow b$），或者经由共同原因与 $a$、$b$ 相关联，形成诸如 $a\leftrightarrow c\leftrightarrow b$ 的模式。正是这一条件使 IC* 算法得以开始为图中的边定向（第 2 步），并添加指向 $c$ 的箭头。也正是这一非传递模式，被用来在定义 2.7.1 中确保 $X$ 不是 $Y$ 的后果、在定义 2.7.2 中确保 $Z$ 不是 $X$ 的后果。在定义 2.7.3 中我们有两个非传递三元组，$(Z_1,X,Y)$ 与 $(X,Y,Z_2)$，从而排除 $X$ 与 $Y$ 之间的直接影响，于是虚假关联便成为它们相依的唯一解释。

对非传递三元组的这种解释，涉及对\emph{结果}变量而非对假定原因的虚拟控制；这类似于操纵论因果观中对零假设的检验（\S 1.3）。例如，人们坚持是雨使草变湿、而不是相反，理由之一是他们可以轻易找到其他使草变湿、且与雨全然无关的手段。转译到我们的链 $a-c-b$，若我们能找到另一条途径（$b$）可以在不影响 $a$ 的情况下潜在地控制 $c$，我们就排除了 $c$ 作为 $a$ 之原因的可能（Pearl 1988a, p.~396）。这一类比当然只是启发性的，因为在观测研究中，我们必须等待自然提供适当的控制，并避免让该控制受到（与 $a$ 的）虚假关联的污染。

\section{非时间性因果与统计时间}

从非时间数据中判定因果影响的方向，会引出若干关于时间与因果解释之关系的有趣哲学问题。例如，定义 2.7.2 或 2.7.4 中赋予箭头 $X\to Y$ 的方向，会不会与可得的时间信息相冲突（比如后来发现 $Y$ 先于 $X$）？既然定义 2.7.4 的理据建立在关于因果关系之统计方面的强烈直觉之上（例如“无某种因果则无相关”），那么这种冲突即便发生，显然也相当罕见。于是问题变成：为什么仅由统计依赖所确定的方向，会与时间之流有任何关系？

在人类话语中，因果解释满足两种期待：时间的与统计的。时间方面体现为“原因应当先于结果”这一理解。统计方面则期待一个完整的因果解释能屏蔽其各个结果（即使结果之间条件独立）；\footnote{这一期待被称为 Reichenbach 的“合取分叉”或“共同原因”判据（Reichenbach 1956; Suppes and Zaniotti 1981; Sober and Barrett 1992），曾受到 Salmon (1984a) 的批评，他表明有些事件虽有资格作为因果解释，却不满足 Reichenbach 的判据。然而 Salmon 的例子涉及不完整的解释，因为它们遗漏了在原因与其各结果之间起中介作用的变量（见 \S 2.9.1）。}不能屏蔽其结果的解释被视为“不完整”，残余的依赖被视为“虚假的”或“未解释的”。这两种期待历经数个世纪的科学观察而毫无冲突地共存，这意味着自然现象的统计必定呈现出某种基本的时间偏向。确实，我们常常遇到这样的现象：对当前状态的知识使未来状态的变量条件独立（例如 \textup{(2.3)} 那样的多元经济时间序列）。然而，我们很少发现相反的现象：对当前状态的知识使过去状态的各分量条件独立。这种时间偏向有什么必然的理由吗？

刻画这一偏向的一种便利方式，是借助\emph{统计时间}的概念。

\begin{definition}[统计时间]
给定经验分布 $P$，$P$ 的\emph{统计时间}是变量的任何一种与至少一个和 $P$ 相容的极小因果结构相一致的排序。
\end{definition}

例如我们看到，标量马尔可夫链过程有许多统计时间：一个与物理时间一致，一个与之相反，还有其他对应于“从某一节点（任意选作根）出发沿马尔可夫链的任一方向”的排序。另一方面，由两条耦合马尔可夫链支配的过程，例如
\begin{equation}
X_t=\alpha X_{t-1}+\beta Y_{t-1}+j_t,\qquad
Y_t=gX_{t-1}+dY_{t-1}+h_t, \tag{2.3}
\end{equation}
却只有一个统计时间——与物理时间一致的那一个。\footnote{这里假定 $j_t$ 与 $h_t$ 是两个相互独立的白噪声时间序列，且 $\alpha=d$、$g=\beta$。}确实，对从这种过程取样的样本运行 IC 算法（同时压制一切时间信息），会很快把 $X_{t-1}$ 与 $Y_{t-1}$ 的各分量识别为 $X_t$ 与 $Y_t$ 的真正原因。这可由定义 2.7.1（其中 $X_{t-2}$ 以 $Z=Y_{t-2}$、$S=\{X_{t-3},Y_{t-3}\}$ 取得 $X_{t-1}$ 之可能原因的资格）与定义 2.7.2（其中 $X_{t-1}$ 以 $Z=X_{t-2}$、$S=\{Y_{t-1}\}$ 取得 $X_t$ 之真正原因的资格）看出。

前面所假定的时间偏向可以表述如下。

\begin{conjecture}[时间偏向]
在大多数自然现象中，物理时间与至少一个统计时间相重合。
\end{conjecture}

Reichenbach (1956) 把与其合取分叉相联系的不对称性归因于热力学第二定律。第二定律能否充分说明刚才描述的时间偏向，是可疑的，因为外部噪声 $j_t$ 与 $h_t$ 的影响使 \textup{(2.3)} 中的过程成为非保守的。\footnote{这一观察得益于 Seth Lloyd。}此外，时间偏向还依赖于语言。例如，把 \textup{(2.3)} 用不同的坐标系表达——比如作线性变换
\[
X'_t=aX_t+bY_t,\qquad Y'_t=cX_t+dY_t,
\]
就有可能使 $(X',Y')$ 表示中的统计时间与物理时间相反；也就是说，$X'_t$ 与 $Y'_t$ 将在给定其\emph{未来}值 $(X'_{t+1},Y'_{t+1})$ 而非过去值时彼此独立。这提示：物理时间与统计时间的一致吻合，是人类选择语言基元所造成的副产品，而不是物理实在的特征。例如，若 $X_t$ 与 $Y_t$ 表示两个相互作用粒子在时刻 $t$ 的位置，$X'_t$ 表示其质心位置、$Y'_t$ 表示其相对距离，那么用 $(X,Y)$ 还是 $(X',Y')$ 坐标系描述粒子运动，（原则上）只是一个选择问题。然而显然，这一选择并非全然任意；它反映出一种偏好：偏好那些使前向扰动（\textup{(2.3)} 中的 $j_t$ 与 $h_t$）彼此正交、而不是使相应的后向扰动（$j'_t$ 与 $h'_t$）彼此正交的坐标系。Pearl and Verma (1991) 推测，这种偏好代表着便于预测未来事件的生存压力，而演化显然把这种能力排在对当前事件作事后解释的能力之前。究竟是这一力量还是别的力量塑造了我们的语言选择，仍有待研究（见 Price 1996 的讨论），而这恰恰使统计—时间的一致性更加引人入胜。

\section{结论}

本章所呈现的理论表明：尽管统计分析并非在一切可以设想的情形中都能把真正的因果关系与虚假共变区分开来，但在许多情形中它能够做到。在模型极小性（和/或稳定性）的假设下，存在某些依赖模式，足以揭示真正的因果关系。这些关系不能被归因于隐藏原因，否则我们就违反了科学方法论的一条基本箴言：以语义方式表述的奥卡姆剃刀。遵守这条箴言，或许可以解释：面对与观测完全相容的种种对立备选，人类为何仍能就因果关系的方向性与非虚假性达成共识。呼应 Cartwright (1989)，我们用一句口号概括自己的主张：“输入无因——输出无因；奥卡姆剃刀入——若干原因出。”

IC 算法——或 Spirtes 等 (1993) 的 TETRAD 程序、或 Cooper and Herskovits (1991)、Heckerman 等 (1994) 的贝叶斯方法——所推断出的因果关系有多可靠？

把这个问题的语境换到视觉知觉上，我们同样可以问：当我们从二维阴影、或从物体在视网膜上投下的二维图像识别三维物体时，我们的预测有多可靠？答案是：并非绝对可靠，但足以分辨树木与房屋，足以在不触摸所见每一物体的前提下作出有用的推断。回到因果推断，我们的问题便成为：典型的学习环境（比如技能习得任务或流行病学研究）中是否有足够多的判别线索，使我们能在原因与结果之间作出可靠区分？若改述为一个逻辑保证，我们可以断然宣称：只要 $a$ 对 $b$ 确实没有因果影响，且所观测分布相对于其底层因果模型是稳定的，IC* 算法就绝不会把箭头 $a\to b$ 标记为真正的。

在实践方面，我们已经表明：模型极小性假设连同“稳定性”假设（无偶然独立），导致了一个有效的算法，用以构造能生成数据的候选因果模型——无论是显式的还是潜在的。1990 年在我们实验室进行的模拟研究表明，含有数十个变量的网络，只需少于 5000 个样本即可由该算法恢复其结构。例如，从 \textup{(2.3)} 所示过程（的二值版本）中取 1000 个样本，每个样本含十对相继的 $X,Y$ 值，就足以恢复其双链结构（以及正确的时间方向）。噪声越大，恢复越快（到某一程度为止）。在把这一建模方案用于真实数据检验时，我们考察了 Sewall Wright 那篇开创性论文《玉米与生猪的相关》（Wright 1925）中所报告的观测。正如所料，玉米价格（$X$）可以被清楚地识别为生猪价格（$Y$）的原因，反之则不然。原因在于存在变量玉米收成（$Z$），它满足定义 2.7.2 的条件（取 $S=\varnothing$）。本章所讨论的原则与算法的若干应用，见 Glymour and Cooper (1999, pp.~441--541)。

自然应当追问：这些新的因果判据如何裨益当前的机器学习与数据挖掘研究？在某种意义上，我们的方法类似于标准的、在假设空间中搜索的机器学习方法（Mitchell 1982），其中每个假设代表一个因果模型。遗憾的是，相似性到此为止。机器学习文献中的主流范式，是把每个假设（或理论、或概念）定义为可观测实例的一个子集；一旦我们观测到该子集的全部外延，假设便被无歧义地确定。因果发现并非如此。即便训练样本穷尽了假设子集（在我们的情形中，这相当于精确观测到 $P$），我们仍面对大量等价的因果理论，每一个都规定了截然不同的一组因果断言。因此，对数据的拟合度不足以成为验证因果理论的标准。传统学习任务试图从一组实例推广到另一组实例，而因果建模的任务则是从一组条件下的行为推广到另一组条件下的行为。因此，因果模型应当用一个挑战其对抗条件变化的稳定性的标准来选择，而这正是科学家试图通过受控实验所达成的事情。在没有此类实验时，人们所能做的最好的事，就是依靠虚拟控制变量，就像自然通过定义 2.7.1—2.7.4 的依赖模式所揭示的那些变量。

\subsection{论极小性、马尔可夫性与稳定性}

从关联推断因果的想法，不可能不受到按传统学说训练的科学家们的质疑。自然，本章所述理论背后的假设——极小性与稳定性——也遭到统计学家与哲学家的攻击。本节包含为这些假设辩护的补充思考。

尽管很少有人质疑极小性原则（质疑它无异于质疑科学归纳），但对我们所定义的极小化对象——即因果模型——却有人提出异议。定义 2.2.2 假定随机项 $u_i$ 相互独立，这一假设赋予每个模型马尔可夫性：在给定其父节点（直接原因）时，每个变量与其非后代独立。这意味着，在 d-分离的其他推论之中，若干通常与 Reichenbach (1956) 共同原因原则相联系的、关于因果与关联的熟悉关系也成立——例如“无因果则无相关”“原因屏蔽其结果”“无超距作用”。

如我们在讨论定义 2.2.2 时所解释的，马尔可夫假设是一种约定，用以区分完整模型与不完整模型。\footnote{关于涉及环与选择偏倚的某些非马尔可夫模型的发现算法，已见 Spirtes 等 (1995) 与 Richardson (1996)。}把马尔可夫假设纳入完整因果模型的定义（定义 2.2.2），再通过潜结构（定义 2.3.2）放松这一假设，我们就声明：我们准备承受这样的损失——无法发现那些不能被描述为潜结构的非马尔可夫因果模型。我并不认为这一损失十分严重，因为这类模型——即使宏观世界中确实存在——作为决策指南的效用也很有限。例如，除了事先逐一列出每一种可以设想的干预的效果之外，很难看出人们如何从这样的模型预测干预的效果。

因此并不奇怪，对马尔可夫假设的批评——最著名的是 Cartwright (1995a, 1997) 与 Lemmer (1993) 的批评——有两点共同特征：
\begin{enumerate}[label=\arabic*.,leftmargin=1.8em]
\item 它们提出宏观的非马尔可夫反例，而这些反例可化归为 Salmon (1984) 所考虑的那类马尔可夫潜结构，即交互分叉；
\item 它们没有提出任何可供预测行动与行动组合之后果的替代性非马尔可夫模型。
\end{enumerate}

交互分叉模型见图 2.6(a)。若中间节点 $d$ 未被观测（或未被命名），人们便容易断定马尔可夫假设被违反了，因为被观测的原因（$a$）并没有屏蔽其结果（$b$ 与 $c$）。图 2.6(b) 的潜结构可以在一切方面模拟图 2.6(a)，二者无论在观测上还是在实验上都无法区分。

\begin{figure}[htbp]
\centering
\begin{tikzpicture}[causal,scale=0.95,every node/.style={transform shape}]
  % (a)
  \begin{scope}
    \node[circle,fill,inner sep=1.3pt] (a) at (-1.6,0) {}; \node[font=\small,left=0.4mm] at (-1.6,0) {$a$};
    \node[lat] (d) at (0,0) {$d$};
    \node[circle,fill,inner sep=1.3pt] (b) at (1.7,0.6) {}; \node[font=\small,right=0.4mm] at (1.7,0.6) {$b$};
    \node[circle,fill,inner sep=1.3pt] (c) at (1.7,-0.6) {}; \node[font=\small,right=0.4mm] at (1.7,-0.6) {$c$};
    \draw[ed] (a)--(d); \draw[ed] (d)--(b); \draw[ed] (d)--(c);
    \node[font=\small] at (0,-1.5) {(a)};
  \end{scope}
  % (b)
  \begin{scope}[xshift=6.4cm]
    \node[circle,fill,inner sep=1.3pt] (a2) at (-1.6,0) {}; \node[font=\small,left=0.4mm] at (-1.6,0) {$a$};
    \node[circle,fill,inner sep=1.3pt] (b2) at (1.7,0.6) {}; \node[font=\small,right=0.4mm] at (1.7,0.6) {$b$};
    \node[circle,fill,inner sep=1.3pt] (c2) at (1.7,-0.6) {}; \node[font=\small,right=0.4mm] at (1.7,-0.6) {$c$};
    \draw[ed] (a2)--(b2); \draw[ed] (a2)--(c2);
    \draw[bid,dashed] (b2) to[bend left=28] (c2);
    \node[font=\small] at (0,-1.5) {(b)};
  \end{scope}
\end{tikzpicture}
\caption{(a) 交互分叉。(b) 与 (a) 等价的潜结构。}
\label{fig:2.6}
\end{figure}

只有量子力学现象才表现出无法归因于潜变量的关联；若有人在宏观世界中发现这类奇特关联，那将被视为科学上的奇迹。尽管如此，马尔可夫条件的批评者仍坚持认为，某些所谓的反例必须建模为 $P(bc\mid a)$ 而非 $\sum_d P(b\mid d,a)P(c\mid d,a)$——也许他们假定，把 $b$ 与 $c$ 之间的依赖留作未解释，会获得某种洞见或一般性。

具有讽刺意味的是，马尔可夫条件之普遍性的最强证据，或许恰恰见于那个称为“概率因果性”的哲学纲领（见 \S 7.5），而 Cartwright 正是其主要倡导者之一。在该纲领中，因果依赖被定义为：在对某组相关因素取条件之后仍然持续的概率依赖（Good 1961; Suppes 1970; Skyrms 1980; Cartwright 1983; Eells 1991）。这一定义立足于这样一个假设：对正确的因素集取条件，可以压制一切虚假关联——而这正等价于马尔可夫条件。概率因果性作为一项活跃的哲学纲领历经四十年而不衰，这一思想史事实恰恰证明：马尔可夫条件的反例相对罕见，且可以通过潜变量加以解释。

现在我来讨论稳定性假设。为稳定性辩护的常见论证（Spirtes et al. 1993）诉诸这样一个事实：在参数可以彼此独立变化的任何概率空间中，参数乘积之间的严格等式具有零勒贝格测度。例如，模型 \textup{(2.2)} 中的等式 $\alpha=-\gamma g$，若考虑参数 $\alpha,-\gamma,g$ 上的任一连续联合密度，其成立的概率为零——除非该密度以某种方式先天地体现了约束 $\alpha=-\gamma g$。与此相反，Freedman (1997) 声称，没有理由假定参数事实上不会被此类等式约束绑定在一起；若如此，所得分布便（依定义 2.4.1）是不稳定的。

确实，因果模型所宣示的条件独立，无非是联合分布上的等式约束。例如链模型 $Y\to X\to Z$ 蕴含等式
\[
r_{YZ}=r_{XZ}\cdot r_{YX},
\]
其中 $r_{XY}$ 是 $X$ 与 $Y$ 之间的相关系数；这一等式约束把三个相关系数永久地绑在一起。那么，究竟是什么使相关系数之间的等式比另一组参数（比如 $\alpha,b,g$）之间的等式享有特权？为什么我们把等式 $r_{YZ}=r_{XZ}\cdot r_{YX}$ 视为“稳定的”，而把 $\alpha=-\gamma g$ 视为“偶然的”？

答案当然在于自主性概念（Aldrich 1989），这一概念处于一切因果概念的核心（见 \S 1.3 与 \S 1.4）。因果模型并不只是通过一组参数编码概率分布的另一种方案。因果模型的独特之处在于：每个变量都通过一种关系（称为“机制”）由另一组变量所决定，而当其他机制受到外部影响时，该关系保持不变。这种不变性意味着机制可以彼此独立地变化，进而意味着当实验条件改变时，结构系数集（例如 \textup{(2.2)} 例子中的 $\alpha,b,g$）——而不是其他类型的参数（如 $r_{YZ},r_{XZ},r_{YX}$）——能够并且将会独立地变化。因此，形如 $\alpha=-\gamma g$ 的等式约束与自主性观念相悖，在自然条件下极少出现。正是人们对稳定性的追求，解释了为什么他们除了设想两个独立原因的共同结果之外，无法举出任何非传递依赖模式的实例（见第 43 页）。任何能令人信服地例示某一依赖模式的故事，都必须在该故事参数取任何数值时都维持这一模式——这正是我们称为“稳定性”的要求。

出于这一理由，有人建议：仅基于关联的因果发现方法，如 IC* 算法或 TETRAD-II 程序所体现的那些，其最大潜力将在条件略有变化的纵向研究中实现——在那里，偶然的独立被破坏，只有结构性的独立得以保留。这一建议假定：在此类变化条件下，模型的参数会受到扰动而其结构保持不变——这是一种微妙的平衡，或许难以验证。尽管如此，考虑到只能依赖受控随机实验这一替代方案，这类纵向研究仍是一个令人兴奋的机会。

\subsection*{与贝叶斯进路的关系}

必须强调，极小性与稳定性原则的要素同样构成贝叶斯进路中因果发现的基础。在该进路中，人们依据结构与参数为一组候选因果网络赋予先验概率，再用贝叶斯法则为给定网络拟合数据的程度打分（Cooper and Herskovits 1991; Heckerman et al. 1999）。随后在可能结构的空间中进行搜索，以寻求后验得分最高的那个（些）结构。基于这一进路的方法具有在小样本条件下表现良好的优点，但在处理隐藏变量时遇到困难。贝叶斯进路的一切实际实现都作出参数独立性假设，这一假设诱导出对参数更少之模型的偏好，因而偏好极小性。同样，参数独立性只有在参数所代表的机制可以彼此自由变化时——即系统是自主的、因而稳定时——才能得到辩护。

\section*{第二版后记}
\addcontentsline{toc}{section}{第二版后记}

卡内基梅隆大学的 TETRAD 小组一直在积极推进因果发现的研究，其成果见 Spirtes 等 (2000)、Robins 等 (2003)、Scheines (2002) 以及 Moneta and Spirtes (2006)。Spirtes、Glymour、Scheines 与 Tillman (2010) 总结了因果发现的当前技术现状。

因果发现在经济学中的应用见 Bessler (2002)、Swanson and Granger (1997) 以及 Demiralp and Hoover (2003)。Gopnik 等 (2004) 运用因果贝叶斯网络来解释儿童如何从观察与行动中获得因果知识（另见 Glymour 2001）。

Hoyer 等 (2006) 与 Shimizu 等 (2005, 2006) 提出了一种发现因果方向性的新方案，它不基于条件独立，而基于函数复合。其想法是：在带有非高斯噪声的线性模型 $X\to Y$ 中，变量 $Y$ 是两个独立噪声项的线性组合。因此，$P(y)$ 是两个非高斯分布的卷积，打个比方说，它比 $P(x)$“更高斯”。这种“更高斯”的关系可以被赋予精确的数值度量，并用于推断某些箭头的方向性。

Tian and Pearl (2001a,b) 发展了又一种因果发现方法，它基于对“冲击”的检测——即环境中自发的局部变化，它们起着“自然的干预”的作用，并揭示出指向这些冲击之后果的因果方向性。

Verma and Pearl (1990) 注意到，两个潜结构可能蕴含相同的条件独立集，却对联合分布施加不同的等式约束。这些约束被称为“休眠独立”，由 Tian and Pearl (2002b) 与 Shpitser and Pearl (2008) 系统地刻画；它们有望为结构学习提供一种强大的新发现工具。

一个名为“因果性工作台”（Causality Workbench）的因果发现算法基准测试计划，已由 Guyon 等 (2008a,b; \url{http://clopinet.com/causality}) 报告。该计划定期组织竞赛，向参赛者提供真实数据或由隐藏因果模型生成的数据，挑战在于预测一组选定干预的结果。
